\documentclass[10pt]{article}

\usepackage{graphicx}
\usepackage{amssymb}
\usepackage{amsmath}
\usepackage{amsthm,nicefrac}
\usepackage{mathtools}
\usepackage{verbatim}
\usepackage{enumitem}
\usepackage{url}

%%%%%%%%%%%%%%%%%%%%%%%%%%%%%%%%%%%%%%%
\setlength{\textwidth}{5.325in}
\setlength{\textheight}{8.25in}
\setlength{\topmargin}{-0.4cm}
\setlength{\evensidemargin}{1.6cm}
\setlength{\oddsidemargin}{1.6cm}
%%%%%%%%%%%%%%%%%%%%%%%%%%%%%%%%%%%%%%%

\newcounter{listacnt}\renewcommand{\thelistacnt}{\alph{listacnt}}
\newenvironment{theoremlist}{\begin{list}{\textnormal{(\thelistacnt)}}%
{\settowidth{\labelwidth}{(b)}%
\setlength{\topsep}{1.5pt plus 0.5pt minus 0.5pt}%
\setlength{\itemsep}{1pt plus 0.2pt minus 0.2pt}%
\setlength{\parsep}{0.1pt plus 0.1pt minus 0.1pt}%
\usecounter{listacnt}}}{\end{list}}


% ==================  SHANE MATH DEFINITIONS  ==================
% ---------------------- bb fonts ---------------------- 
\newcommand{\rr}{\mathbb{R}}
\newcommand{\qq}{\mathbb{Q}}
\newcommand{\cc}{\mathbb{C}}
\newcommand{\zz}{\mathbb{Z}}
\newcommand{\nn}{\mathbb{N}}
\newcommand{\dd}{\mathbb{D}}
\newcommand{\pp}{\mathbb{P}}
% ----------------------- cal fonts ----------------------
\newcommand{\cP}{\mathcal{P}}
\newcommand{\cT}{\mathcal{T}}
\newcommand{\cX}{\mathcal{X}}
\newcommand{\cQ}{\mathcal{Q}}
\newcommand{\cF}{\mathcal{F}}
\newcommand{\cU}{\mathcal{U}}
\newcommand{\bigo}{\mathcal{O}}
% ---------------------- interval arithmetic ---------------------- 
\newcommand{\intrr}{\mathbb{IR}}
\newcommand{\intcc}{\mathbb{IC}}
\newcommand{\intval}[1]{\tilde{#1}}
\DeclareMathOperator{\rad}{rad}
% ---------------------- variable flair ---------------------- 
\newcommand{\overbar}[1]{\mkern 1.5mu\overline{\mkern-1.5mu#1\mkern-1.5mu}\mkern 1.5mu}
% ---------------------- operator names ---------------------- 
%\DeclareMathOperator{\ker}{ker}
\DeclareMathOperator{\coker}{coker}
\DeclareMathOperator{\image}{image}
\DeclareMathOperator{\sign}{sign}
\DeclareMathOperator{\inspan}{span}
% ---------------------- delimeters ---------------------- 
\newcommand{\setof}[1]{\left\{#1\right\}}
\newcommand{\pardiff}[2]{\frac{\partial #1}{\partial #2}}
\newcommand{\abs}[1]{\left| #1 \right|}
\DeclarePairedDelimiter{\norm}{\lVert}{\rVert}
\DeclarePairedDelimiter{\dotp}{\langle}{\rangle}
%\newcommand{\dotp}[1]{\langle #1 \rangle}
\newcommand{\paren}[1]{\left( #1 \right)}
% ---------------------- misc ---------------------- 
\renewcommand{\Re}{\text{Re}}
\renewcommand{\Im}{\text{Im}}
\renewcommand{\emptyset}{\O}
\newcommand{\ptlt}{\hskip 5pt < \hskip-7pt \cdot} % pointwise less than
\newcommand{\ptgt}{\hskip 5pt \cdot \hskip-7pt >} % pointwise greater than
\newcommand{\bydef}{\,\stackrel{\mbox{\tiny\textnormal{\raisebox{0ex}[0ex][0ex]{def}}}}{=}\,}
% ---------------------- theorem environments ---------------------- 
%\numberwithin{equation}{section} UNCOMMENT TO CHANGE EQUATION NUMBERING 

% amsmath theorem style (italics)
\newtheorem{theorem}{Theorem}%[section]
\newtheorem{corollary}[theorem]{Corollary}
\newtheorem{lemma}[theorem]{Lemma}
\newtheorem{proposition}[theorem]{Proposition}

% amsmath  definition style (no italics)
\theoremstyle{definition}
\newtheorem{definition}{Definition}
\newtheorem{example}{Example}

% amsmath  remark style (no italics, no space)
\theoremstyle{remark}
\newtheorem{remark}{Remark}


\UseRawInputEncoding % Forces compatiblitity with new latex UTF encoding
\begin{document}
\title{Ideal Domain Notes}%
\author{Shane Kepley}
\date{}
\maketitle


The standard coordinate potential is given by
\[
\cU_0 = \mu \paren{\frac{1}{2}r_2^2 + \frac{1}{r_2}} + (1-\mu) \paren{\frac{1}{2}r_1^2 + \frac{1}{r_1}}
\]
where
\[
r_1 = \sqrt{\paren{x_0 - \mu}^2 + y_0^2} \qquad r_2 = \sqrt{\paren{x_0 + 1 - \mu}^2 + y_0^2}
\]
The corresponding energy is
\[
H_0 = 2 \cU_0 - p_0^2 - q_0^2.
\]
We want to compute $\norm{D_{u_0} H_0}^2$ where $u_0 = (x_0, p_0, y_0, q_0)$ is the state in standard coordinates. Start with partials of $\cU_0$ with respect to $x_0,y_0$ given by 
\[
\pardiff{\cU_0}{x_0} = \mu \paren{x_0 + 1 - \mu} \paren{1 - r_2^{-\nicefrac{3}{2}}} + (1-\mu)(x_0 - \mu) \paren{1 - r_1^{-\nicefrac{3}{2}}} 
\]
\[
\pardiff{\cU_0}{y_0} = \mu y_0 \paren{1 - r_2^{-\nicefrac{3}{2}}} + (1-\mu) y_0 \paren{1 - r_1^{-\nicefrac{3}{2}}}
\]
So now, we compute $D_{u_0} H_0$ which is 
\begin{align*}
	D_{u_0} H_0 = 
	2
	\begin{pmatrix}
	\pardiff{\cU_0}{x_0} \\
	-p_0 \\
	\pardiff{\cU_0}{y_0} \\
	-q_0
	\end{pmatrix}
	=
	2
	\begin{pmatrix}
	\mu \paren{x_0 + 1 - \mu} \paren{1 - r_2^{-\nicefrac{3}{2}}} + (1-\mu)(x_0 - \mu) \paren{1 - r_1^{-\nicefrac{3}{2}}}  \\
	-p_0 \\
	 \mu y_0 \paren{1 - r_2^{-\nicefrac{3}{2}}} + (1-\mu) y_0 \paren{1 - r_1^{-\nicefrac{3}{2}}} \\
	 -q_0
	\end{pmatrix}
\end{align*}
Therefore, we get 
\begin{align*}
	\frac{1}{4} \norm{D_{u_0} H_0}^2 & = 
	p_0^2 + q_0^2 + \mu^2 \paren{1 - r_2^{-\nicefrac{3}{2}}}^2 \paren{(x_0 + 1 - \mu)^2 + y_0^2} + \paren{1-\mu}^2  \paren{1 - r_1^{-\nicefrac{3}{2}}}^2 \paren{(x_0 - \mu)^2 + y_0^2} \\
	& = p_0^2 + q_0^2 + \mu^2 \paren{1 - 2 r_2^{\nicefrac{1}{2}} + r_2^{-1}}  + \paren{1 - \mu}^2  \paren{1 - 2r_1^{\nicefrac{1}{2}} + r_1^{-1}}.
\end{align*}
This quantity is an estimate for the energy drift occurring when making an infinitesimal error during integration. 

We want to compare this with the quantity $\norm {D_{u_1} H_0}^2$ where $u_1 = (x_1, p_1, y_1, q_1)$ which estimates the energy loss (with respect to the standard coordinates) due to integrating a point by first regularizing it to $u_1$ coordinates. By the chain rule we have for $u_1^{(j)}$ for $j = 1,2,3,4$
\[
\pardiff{H_0}{u_1^{(j)}} = D_{u_0} H_0 \cdot \frac{d u_0}{du_1^{(j)}}
\]
So we obtain
\[
D_{u_1} H_0 = D_{u_0} H_0 \cdot D_{u_1} u_0
\] 
where $D_{u_1} u_0$ is the $4 \times 4$ derivative of $u_1$ with respect to $u_0$ and $D_{u_0} H_0$ is the vector previously computed. So we compute
\[
D_{u_1} u_0 = 
\begin{pmatrix}
2 x_1 & 0 & -2 y_1 & 0 \\
\pardiff{p_0}{x_1} & \frac{x_1}{2(x_1^2 + y_1^2)} & \pardiff{p_0}{y_1} & - \frac{y_1}{2(x_1^2 + y_1^2)} \\
2 y_1 & 0 & 2 x_1 & 0 \\
\pardiff{q_0}{x_1} & \frac{y_1}{2(x_1^2 + y_1^2)} & \pardiff{q_0}{y_1} & - \frac{x_1}{2(x_1^2 + y_1^2)} \\
\end{pmatrix}
\]
where the partial derivatives for $p_0,q_0$ are given by 
\begin{align*}
	\pardiff{p_0}{x_1} & =  \frac{p_1 (y_1^2 - x_1^2)}{2 (x_1^2 + y_1^2)^2} \\
	\pardiff{p_0}{y_1} & =  \frac{q_1 (y_1^2 - x_1^2)}{2 (x_1^2 + y_1^2)^2} \\
	\pardiff{q_0}{x_1} & =  \frac{q_1 (y_1^2 - x_1^2)}{2 (x_1^2 + y_1^2)^2} \\
	\pardiff{q_0}{y_1} & =  -\frac{p_1 (y_1^2 - x_1^2)}{2 (x_1^2 + y_1^2)^2} \\
\end{align*}
Note the identities
\[y_1^2 - x_1^2 = \mu - x_0 \qquad x_1^2 + y_1^2 = \abs{\sqrt{x_0 - \mu + i y_0}}^2 = (x_0 - \mu)^2 + y_0^2 = r_2^2
\]
and also $x_1 = \frac{y_0}{\mu - x_0 + r_1}$ and $y_1 = \frac{\mu - x_0 + r_1}{2}$.
So we can rewrite $D_{u_1} u_0$ back in terms of $u_0$ to measure the $H_0$ energy drift resulting from integrating a point $u_0$ by first mapping into the $f_1$ coordinates. 
\[
D_{u_1} u_0 \bigg|_{u_1 = \cF_1(u_0)} = 
\begin{pmatrix}
2 \Re(\sqrt{x_0 - \mu + i y_0}) & 0 & -2 \Im(\sqrt{x_0 - \mu + i y_0}) & 0 \\
BLAH & blah & blah & blah \\
\end{pmatrix}
\]
\end{document}

